\chapter{Modellazione di sistemi}

In questo capitolo, ci concentreremo sulla modellazione di sistemi software, che rappresenta un passo fondamentale per la verifica formale.

Come anticipato nel capitolo precedente, la verifica formale si base sul \keyword{model checking}, il quale ovviamente necessita di un modello formale del software da verificare.
Utilizzare questa tecnica di verifica formale, ovviamente, cambia completamente il processo di sviluppo software rispetto a quello visto nella \Cref{fig:testing_cycle}, portandolo a qualcosa di più simile a quello in \Cref{fig:formal_verification_cycle}.

\begin{figure}[H]
  \centering
  \includegraphics[width=\textwidth]{_images/formal_verification_cycle}
  \caption{Ciclo di sviluppo software con verifica formale}\label{fig:formal_verification_cycle}
\end{figure}

Come illustrato in \Cref{fig:formal_verification_cycle}, il processo si articola in due fasi parallele che convergono nel Model Checking.

\paragraph{Formalizzazione dei Requisiti:}
Partendo dai requisiti del sistema, intesi come le proprietà che devono essere soddisfatte, attraverso un processo di \keyword{formalizzazione} si producono \keyword{specifiche di proprietà} (property specification) espresse come formule in un linguaggio logico formale.

Questo passaggio è tutt'altro che banale: i requisiti sono tipicamente espressi in linguaggio naturale o semi-formale, mentre i metodi formali richiedono necessariamente linguaggi formali. La scelta del linguaggio dipende dal tipo di requisiti e dalle proprietà che si vogliono esprimere.

\paragraph{Modellazione del Sistema:}
Parallelamente, dalla descrizione del sistema, spesso espressa tramite i linguaggi di progettazione dell'ingegneria del software, attraverso un processo di \keyword{modellazione} (modeling) si costruisce un \keyword{modello formale del sistema} (system model).

\paragraph{Coerenza tra Modello e Specifiche:}
È fondamentale che il modello del sistema sia \keyword{coerente} con la semantica del linguaggio di specifica delle proprietà. Questo significa che il modello deve contenere sufficiente \qi{ricchezza matematica} per poter interpretare le formule: deve fornire abbastanza oggetti matematici per dare significato ai simboli logici utilizzati nelle specifiche.

Il modello può essere più ricco di quanto strettamente necessario per interpretare le formule, ma non può essere meno ricco, altrimenti la verifica non sarebbe possibile.

\paragraph{Model Checking e Analisi dei Risultati:}
Una volta completate le fasi di formalizzazione e modellazione, il \keyword{Model Checking} verifica se il modello del sistema soddisfa le proprietà specificate (\(\mathcal{M} \models \varphi\)). Il processo può produrre due esiti:
\begin{itemize}
  \item \textbf{Satisfied}: la proprietà è soddisfatta dal modello;
  \item \textbf{Violated}: la proprietà è violata, e viene prodotto un \keyword{controesempio} (counterexample) che identifica una traccia di esecuzione che viola la proprietà.
\end{itemize}

Il controesempio può essere analizzato tramite simulazione per identificare la causa dell'errore, che può risiedere nel sistema stesso o nella sua modellazione.

Questo problema ovviamente è legato alla ineliminabile (in linea di principio) discordanza o \q{gap} tra l'astrazione eil sistema.
Il modello potrebbe contenere delle esecuzioni in più di quelle che sono concrete.

Non è detto che però non sia una cosa desiderabile che ci sia questo gap, o quantomeno potrebbe essere impossibile di fatto evitarlo per questioni computazionali.
Questo poiché, avere un modello che ammetta più comportamenti di quelli concreti, non porta a una maggiore complicazione della descrizione, anzi, tipicamente porta una riduzione della complessità della descrizione.
Infatti, meno informazione sono inserite nel modello, più comportamenti sono ammessi, ma maggiore è l'informazione, più vincoli da controllare stanno venendo inseriti nell'insieme delle computazioni.
Ciò che è fondamentale è quindi fare in modo che il modello contenga certamente tutti i comportamenti del sistema concreto, ma potenzialmente anche di più.

In un certo senso è possibile trarre un'analogia intuitiva col concetto di astrazione: il modello è come
un'astrazione più o meno lontana dalla realtà, dalla verità. Quindi più l'astrazione
è lontana dalla realtà, più comportamenti ammette e più semplice è la sua
descrizione.

Ovviamente, questa necessità di bilanciare livello di dettaglio e complessità computazoinale del problema ha fatto nascere delle tecniche di \textit{abstraction refinement}.
Queste partono da una descrizione astratta e, in base ai testimoni spuri che naturalmente vengono prodotti, richiedono dei raffinamenti mirati ad escludere quelle computazioni spurie.

\section{Modellazione}

Per poter costruire modelli di sistemi software, è necessatio in primo luogo definire un linguaggio per la descrizione di tali modelli.

Un linguaggi di modellazione di base è quello dei \keyword{transition system}.

\begin{definition}{Sistemi a Transizioni}
  Un \keyword{sistema a transizioni} o \keyword{transition system} è una terna \((S,S_0,R)\) dove:
  \begin{itemize}
    \item \(S\) è un insieme di stati;
    \item \(S_0 \subseteq S\) è un insieme di stati iniziali;
    \item \(R \subseteq S \times S\) è una endo-relazione di transizione totale a sinistra (o seriale), ovvero tale che:
          \[\forall s \in S ,\exists s' \in S : (s,s') \in R\]
  \end{itemize}
  In poche parole non esistono stati senza transizioni uscenti, o stati terminali.
\end{definition}
La scelta di non ammettere stati terminali è motivata dal fatto che, in questo modo, è possibile modellare anche sistemi che non terminano mai, come ad esempio sistemi operativi o server.
La terminazione è ancora modellabile anche in questi modelli come proprietà del sistema.

È importante notare che \((S,R)\) definiscono un grafo orientato, quindi è ben definito il concetto di cammino o percorso, e dunque di raggiungibilità.

\begin{note}{}
  Una notazione alternativa per la relazione di transizione è quella di rappresentarla con una freccia, ovvero \(s \xrightarrow{} s'\) per indicare che \((s,s')\in R\).
\end{note}

\subsection{Sistemi Chiusi e Sistemi Aperti}

Il modello dei sistemi a transizioni rappresenta la struttura più elementare e generale per la modellazione di sistemi computazionali. È possibile osservare come questo modello sia strettamente correlato al concetto di automa, costituendone essenzialmente il nucleo centrale.

\paragraph{Relazione con gli Automi a Stati Finiti:}
Un automa a stati finiti classico per linguaggi regolari possiede tipicamente due componenti aggiuntive rispetto a un sistema a transizioni:

\begin{itemize}
  \item un \keyword{alfabeto} di simboli che etichettano le transizioni;
  \item una \keyword{condizione di accettazione}, ovvero un insieme di stati finali che determinano l'accettazione di una stringa.
\end{itemize}

Un sistema a transizioni può essere visto come un automa in cui l'intero insieme \(S\) è accettante e in cui esiste un unico simbolo (implicito) nell'alfabeto.

Se si introducesse esplicitamente un alfabeto, la relazione di transizione \(R\) dovrebbe diventare una relazione ternaria \(R \subseteq S \times \Sigma \times S\), dove \(\Sigma\) rappresenta l'alfabeto. Tuttavia, quando l'alfabeto contiene un unico simbolo, la presenza di tale simbolo non modifica l'insieme delle coppie di stati raggiungibili, rendendo quindi superflua la sua esplicitazione. La situazione sarebbe differente con più simboli: da uno stato \(s\) con un simbolo \(\sigma_1\) si potrebbe raggiungere uno stato \(s'\), mentre con un simbolo \(\sigma_2\) tale transizione potrebbe non essere ammessa.

\paragraph{Generalità del Modello:}
I sistemi a transizioni costituiscono un modello sufficientemente generale per rappresentare sistemi computazionali. Gli automi tradizionali, come studiati nella teoria dei linguaggi formali, sono primariamente \keyword{riconoscitori di linguaggi}, progettati per accettare o rifiutare stringhe.
I sistemi computazionali, analogamente alle macchine di Turing, hanno uno scopo più ampio: non solo riconoscere, ma anche \keyword{calcolare}.

In una macchina di Turing, l'alfabeto non rappresenta necessariamente simboli letti dall'esterno: il sistema riceve un input iniziale e poi evolve autonomamente operando sul proprio nastro.
Similmente, in un sistema a transizioni non è necessario il concetto di stimolo o segnale esterno. Negli automi, invece, l'alfabeto rappresenta intuitivamente stimoli esterni che guidano l'evoluzione del sistema.

I sistemi computazionali possono inoltre essere classificati in base alla loro interazione con l'ambiente esterno.

\begin{definition}{Sistemi Chiusi}
  I \keyword{sistemi chiusi} sono sistemi in cui eventuali input sono codificati all'interno degli stati iniziali e il sistema evolve autonomamente in base alla propria logica interna, senza ricevere stimoli esterni durante l'esecuzione.
\end{definition}

\begin{definition}{Sistemi Aperti}
  I \keyword{sistemi aperti} sono sistemi che interagiscono continuamente con l'ambiente esterno, combinando decisioni interne (basate sul proprio stato) con stimoli esterni forniti dall'ambiente durante l'esecuzione.
\end{definition}

\paragraph{Non-Determinismo e Scelte:}
Nei sistemi aperti emergono due tipologie distinte di scelte:
\begin{itemize}
  \item \textbf{Scelte interne}: rappresentano il \keyword{non-determinismo} classico, dove il sistema può evolvere in modi diversi a partire dallo stesso stato, indipendentemente da stimoli esterni;
  \item \textbf{Scelte esterne}: rappresentano la selezione degli stimoli forniti dall'ambiente al sistema.
\end{itemize}

Nel contesto degli automi non deterministici, dato un certo simbolo esterno, l'automa può evolvere in stati diversi: il sistema sceglie internamente come reagire allo stimolo. Quando la scelta del simbolo stesso non dipende dal sistema, ma dall'ambiente, entrambe le tipologie di scelte devono essere considerate.

Per modellare sistemi aperti è necessario quindi prevedere sia le scelte interne del sistema (come evolve da un dato stato), sia le scelte esterne dell'ambiente (quali stimoli fornisce). Questo porta a una interpretazione in termini di \keyword{teoria dei giochi}: un sistema aperto può essere visto come un gioco tra due giocatori avversari, dove il sistema compie le proprie scelte interne e l'ambiente (visto come avversario) sceglie quali stimoli fornire.

\begin{note}{}
  In questo corso ci concentreremo esclusivamente su \keyword{sistemi chiusi}, che possono essere interpretati come giochi a un singolo giocatore, in cui il sistema è l'unico responsabile della risoluzione del proprio non-determinismo. I sistemi aperti, pur essendo un argomento di grande interesse teorico e intimamente legato alla teoria dei giochi, non saranno trattati per vincoli di tempo.
\end{note}

\subsection{Computazioni e Percorsi}

All'interno di un sistema a transizioni è possibile definire formalmente il concetto di computazione.

Come detto precedentemente ogni sistema a transizioni definisce un grafo orientato, i cui percorsi rappresentano esattamente le possibili evoluzioni del sistema, ovvero le sue computazioni.

\begin{definition}{Percorso Finito}
  Un \keyword{percorso finito} da uno stato \(s\) è una sequenza di stati \(s_1, s_2, \ldots, s_n\) tale che \(s = s_1\) e \(s_i \xrightarrow{} s_{i+1}\) per ogni \(i = 1, \ldots, n-1\).
\end{definition}

Poiché la relazione di transizione \(R\) è totale (seriale), ogni percorso massimale è necessariamente infinito. Questo porta naturalmente a distinguere due tipologie di computazioni:
\begin{itemize}
  \item \textbf{Computazioni finite}: corrispondono a percorsi finiti non massimali;
  \item \textbf{Computazioni infinite}: corrispondono a percorsi infiniti massimali.
\end{itemize}

\subsection{Sistemi a Transizioni Etichettati}

Il modello base dei sistemi a transizioni può essere esteso introducendo un insieme finito di \keyword{azioni} che etichettano le transizioni, analogamente agli automi a stati finiti.

\begin{definition}{Labeled Transition System}
  Un \keyword{sistema a transizioni etichettato} (Labeled Transition System, LTS) è una quadrupla \(TS=(S, Act, S_0, R)\) dove:
  \begin{itemize}
    \item \(S\) è un insieme di stati;
    \item \(Act\) è un insieme finito di azioni;
    \item \(S_0 \subseteq S\) è un insieme di stati iniziali;
    \item \(R \subseteq S \times Act \times S\) è una relazione di transizione ternaria.
  \end{itemize}
  La relazione \(R\) indica quale stato può essere raggiunto a partire da un altro stato eseguendo una determinata azione \(a \in Act\). Si scrive \(s \xrightarrow{a} s'\) per indicare che \((s,a,s') \in R\).
\end{definition}

\begin{example}{Vending Machine}
  Consideriamo un esempio concreto di sistema a transizioni etichettato: una macchina distributrice automatica modellata dal LTS in figura.

  \begin{figure}[H]
    \centering
    \includegraphics[width=\textwidth]{_images/vending_machine_example}
    \caption{Vending Machine}\label{fig:vending_machine_example}
  \end{figure}

\end{example}
\subsection{Esempio: Vending Machine}

Questo sistema, sebbene nella realtà debba interagire con l'ambiente esterno, verrà modellato come un sistema chiuso interpretando tutte le azioni come interne. Le azioni del sistema includono:
\begin{itemize}
  \item \textit{coin}: inserimento di una moneta;
  \item \textit{tea}, \textit{coffee}: selezione della bevanda;
  \item \textit{c-serve}, \textit{t-serve}: erogazione rispettivamente di caffè e tè;
  \item \textit{coin-return}: restituzione della moneta.
\end{itemize}

Il comportamento del sistema è il seguente: dallo stato iniziale, dopo l'azione di inserimento della moneta (\textit{coin}), in base alla presenza o assenza della bevanda richiesta, il sistema può scegliere di erogarla o di transire a uno stato di errore dal quale viene eseguito un recovery restituendo la moneta all'utente (\textit{coin-return}).

\paragraph{Non-Determinismo e Astrazione:}
Il modello presenta un non-determinismo apparente: la macchina sembra poter decidere autonomamente se la bevanda è disponibile o meno. Questo non-determinismo è una conseguenza del livello di astrazione scelto, in cui non è stata specificata l'informazione necessaria per discriminare tra le diverse azioni possibili.

In un'implementazione concreta, si utilizzerebbero contatori per tracciare le quantità disponibili di caffè e tè. La descrizione attuale, essendo non deterministica, ammette computazioni spurie. Per eliminare tale non-determinismo, sarebbe necessario arricchire il modello con l'informazione sulla disponibilità delle bevande, introducendo una \keyword{memoria} esplicita nel sistema.

Negli automi a stati finiti, la memoria è codificata negli stati stessi. Per rappresentare esplicitamente le quantità di caffè e tè disponibili, sarebbe necessario moltiplicare il numero di stati per tutte le possibili combinazioni dei valori dei due contatori, portando a un'esplosione combinatoria della dimensione del modello.

Questo esempio illustra un principio fondamentale della modellazione: se le proprietà da verificare non dipendono da determinati aspetti del sistema, tali aspetti possono essere eliminati tramite \keyword{astrazione}, riducendo significativamente la complessità computazionale dell'analisi.

\subsection{Stati di Deadlock}

Nella definizione generale dei sistemi a transizioni, potrebbero esistere stati senza transizioni uscenti, denominati \keyword{stati di deadlock}, in cui la relazione \(R\) non è totale. Nei sistemi concorrenti, gli stati di deadlock possono emergere naturalmente e rappresentare situazioni problematiche del sistema.

Per semplicità, in questo corso assumeremo generalmente che la relazione di transizione sia totale, evitando così complessità addizionali nella semantica del modello.

\section{Strutture di Kripke e Proposizioni Atomiche}

Per poter verificare proprietà sui sistemi a transizioni, è necessario arricchire il modello con informazioni che permettano di esprimere e valutare tali proprietà.

\subsection{Proposizioni Atomiche}

Invece di introdurre variabili esplicite nel sistema di base, si definisce un insieme di \keyword{proposizioni atomiche} (\(AP\)) che rappresentano affermazioni elementari sullo stato del sistema. Ad esempio, una proposizione atomica potrebbe essere \qi{il valore di \(x\) è 5}, denotata come \(x=5\).

Per ogni proposizione atomica è necessario specificare in quali stati essa è vera e in quali è falsa. Questo porta all'introduzione di una \keyword{funzione di etichettamento} \(L\) che associa ad ogni stato \(s\) l'insieme delle proposizioni atomiche vere in quello stato:
\[L: S \rightarrow 2^{AP}\]

\subsection{Strutture di Kripke}

Quando si arricchisce un sistema a transizioni con la funzione di etichettamento \(L\), si ottiene un oggetto che in letteratura prende il nome di \keyword{Struttura di Kripke} (Kripke Structure).

\begin{definition}{Struttura di Kripke}
  Una \keyword{Struttura di Kripke} è una quadrupla \(M=(S, S_0, R, L)\) dove:
  \begin{itemize}
    \item \((S, S_0, R)\) è un sistema a transizioni;
    \item \(L: S \rightarrow 2^{AP}\) è una funzione di etichettamento che associa ad ogni stato l'insieme delle proposizioni atomiche vere in quello stato.
  \end{itemize}
\end{definition}

Le strutture di Kripke devono il loro nome a Saul Kripke, logico che ha definito la semantica delle logiche modali. Le logiche temporali utilizzate in questo corso sono infatti logiche modali, e le strutture di Kripke rappresentano il ponte tra i sistemi a transizioni e i modelli logici.

Le proposizioni atomiche e la funzione di etichettamento \(L\) convertono un sistema a transizioni in un \keyword{modello logico}. Rispetto a questo modello è possibile interpretare e verificare formule espresse in linguaggi logici formali, come verrà discusso nei capitoli successivi.

In una struttura di Kripke, la funzione \(L\) associa a ogni stato un modello proposizionale che rappresenta le proprietà atomiche vere in quello stato.
Possiamo quindi definire un concetto di soddisfacibilità locale, ovvero in uno stato.


\section{Program Graph}
Non tutti i sistemi possono essere mappati in grafi, infatti per sistemi complessi è necessario decomporre il sistema in \q{building blocks} più semplici.
I building blocks \q{atomici} sono i \keyword{Program Graph}, che sono definiti da un insieme di variabili \(\Var = \set{v_1,\ldots,v_k}\) con dominio \(D = \bigcup_{1,\ldots,k} D_i\), con \(D_i\) dominio di \(v_i\).

\begin{note}{}
  Il dominio effettivo delle variabili non è particolarmente importante poiché può essere sempre rappresentato da un insieme numerabile come \(\mathbb{N}\) o \(\mathbb{Z}\).
\end{note}

Le variabili \(\Var\) rappresentano il concetto di memoria, ciò ci permette di modellare sistemi in cui il comportamento dipende dai dati.

I Program Graph invece rappresentano il concetto di un programma, infatti insieme a
i Labeled Transition System, permettono di descrivere come il sistema opera sui dati
durante il flusso di controllo.

\begin{definition}{Program Graph}
  Un \keyword{Program Graph} è una sestupla \((\Loc,\Act,\Eff,\hookrightarrow , \Loc_0, g_0)\) dove:
  \begin{itemize}
    \item \(\Loc\) è un insieme di locazioni o stati di controllo, (dominio del program counter, ovvero l'insieme del numero di linee di codice in un programma);
    \item \(\Act\) l'insieme delle azioni che modificano lo stato del programma;
    \item \(\Eff : \Act\times \Eval(\Var) \to \Eval(\Var)\) è una funzione;
    \item \(\hookrightarrow \subseteq \Loc \times \Cond(\Var) \times \Act \times \Loc\) è la relazione di transizione (che ci permetterà di descrivere il flusso di controllo);
    \item \(\Loc_0 \subseteq \Loc\) è l'insieme delle locazioni iniziali;
    \item \(g_0 \in \Eval(\Var)\) è una condizione che le variabili devono soddisfare all'inizio dell'esecuzione.
  \end{itemize}
\end{definition}

\paragraph{Relazione di transizione (\(\hookrightarrow\)):} intuitivamente collega \q{temporalmente} locazioni o righe del programma. Nello specifico è una relazione che collega due locazioni tramite una condizione sulle variabili, detta anche guardia e un'azione.
La guardia, ovvero la verifica di una certa proprietà sulle variabili, mi dice se è possibile eseguire la transizione, e l'azione mi dice come modificare le variabili quando eseguo la transizione.

\paragraph{La funzione \(\Eff\):} descrive la semantica delle azioni. Prende in input un'azione ed associa ad essa una trasformazione del valore delle variabili prima e dopo l'azione. In altre parole fissata un'azione collega \q{temporalmente} stati delle variabili.


\subsection{Valutazione delle variabili}

Per definire la funzione \(\Eff\) è necessario introdurre un concetto di valutazione delle variabili, e per farlo usiamo \(\Eval(\Var)\), un insieme di funzioni che mappano ogni variabile al suo valore, rendendo \(\Eval(\Var)\) l'insieme di tutte le possibili configurazioni delle variabili.


\subsection{Espressioni}
\(\Act\) e \(\Cond\) sono definite sulla base del concetto di espressioni \(\Expr(\Var\cup D)\), che ci permette di descrivere sia variabili \(x\in\Var\) sia valori costanti del dominio \(x\in D\). Le espressioni dipendono dai domini delle variabili, se prendiamo ad esempio le espressioni aritmetiche abbiamo \(v+1\).


\subsection{Condizioni}
Le condizioni sono dunque valutazioni tra espressioni del tipo \(\expr_1 \cdot \expr_2\), con \(\cdot\) dipendende dal dominio (per aritmetica \(=, \neq, <, > ecc\)), eventualmente anche con operatori logici.

Con gli operatori logici, infatti possiamo creare combinazioni di condizioni, quindi
per esempio possiamo combinarle in congiunzione, disgiunzione, implicazione, negazione, ecc..
Tipicamente utilizzeremo la congiunzione, mentre la negazione verrà applicata internamente alle espressioni atomiche di cui ogni condizione è composta.

\subsection{Azioni}
Le azioni su \(\Var\) sono l'insieme degli assegnamenti, ovvero espressioni del tipo \(v := exp\), dove \(v\in \Var\) e \(exp\in \Expr(\Var\cup D)\).


\subsection{Esempio: Vending Machine come Program Graph}

\begin{example}{Vending Machine come Program Graph}
  Se prendiamo l'esempio \Cref{fig:vending_machine_example} e lo generalizziamo per ottenere un Program Graph, allora abbiamo la seguente figura.

  \begin{figure}[H]
    \centering
    \includegraphics[width=\textwidth]{_images/vending_machine_as_pg}
    \caption{Vending Machine As Program Graph}\label{fig:vending_machine_as_pg}
  \end{figure}

\end{example}


Dove, le variabili \(c\) e \(t\) sono rispettivamente la quantità di caffè a disposizione e la quantità di tè.

Alcune azioni, in particolare quelle non aritmetiche come \(Coin, Coffee\) ecc..\ sono azioni nulle, ovvero non alterano lo stato, ma hanno un impatto sull'evoluzione del sistema da uno stato all'altro.


\section{Compilazione di un Program Graph in un Transition System}

Possiamo notare che i grafi sottostanti a questo \Cref{fig:vending_machine_as_pg} e all'esempio precedente \Cref{fig:vending_machine_example} sono isomorfi, ma i due ogetti sono diversissimi, poiché il secondo sta descrivendo un LTS con stati infiniti (essendo il dominio delle variabili infinito).

È possibile però \q{compilare} un oggetto del genere in una Kripke Structure, ma è necessaria la funzione \(\Eff\) per definire la semantica delle azioni.
Questa compilazione o traduzione, di fatto \q{da semantica} a questa struttura.


Per poter effettuare questa traduzione dobbiamo conoscere le seguenti cose:
\begin{enumerate}
  \item Lo spazio degli stati dell'LTS indotto;
  \item L'insieme delle azioni;
  \item La relazione di transizione dell'LTS indotto;
  \item Gli stati iniziali;
  \item L'etichettatura degli stati.
\end{enumerate}

\subsection{Spazio degli stati}
Per lo spazio degli stati questo equivale ad \(\Eval(\Var)\) (variabili esplicite) e \(\Loc\) (variabili implicite, ovvero il program counter).

\begin{definition}{Valutazione in \(\Eval\)}
  Formalmente una valutazione in \(\Eval\) è una funzione \(\eta: \Var \to D\).
  Inoltre se l'insieme \(\Var\) è finito, una valutazione può essere rappresentata con una tupla lunga col numero di variabile, che accoppiata con una locazione \(l\) rappresenta uno stato del sistema, e dunque un nodo del grafo orientato sottostante.
\end{definition}

Quindi nella Kripke Structure sottostante, \(S = \Loc \times \Eval(\Var)\).

\subsection{Definizione della semantica delle azioni}
Col concetto di \(\Eval(\Var)\) definito possiamo adesso definire \(\Eff\), che intuitivamente ci dice come evolveranno i dati nelle transizioni tra stati.

\begin{definition}{\(\Eff\) non ben definita}
  Data \(\eta \in \Eval(\Var)\) e un azione della forma \(a \overset{\text{def}}{=} v := \expr\), con \(\expr \in \Var \cup D\) allora l'effetto dell'azione è
  \[\Eff(a,\eta)=\eta[v\leftarrow \expr]\] con
  \[\eta[v\leftarrow \expr] = \eta'\]
  dove
  \[\eta' (x)= \begin{cases}
      \expr   & \text{se } x = v  \\
      \eta(x) & \text{altrimenti}
    \end{cases}\]
\end{definition}


Questa definizione non è ben definita perché ci vuole un concetto di valutazione dell'espressione \(\expr\) in \(\Eval(\Var)\). In altre parole un \(\Val(\expr, \eta) \in D\), da definire per induzione strutturale sulla struttura delle espressioni.
I casi base sono costanti e le variabili (valutate con \(\eta\)). Il resto dipende dalla semantica degli operatori.

Quindi, se l'espressione è \(v + 1\), la valutazione di \(v + 1\) rispetto a una certa \(\eta\) sarà il numero corrispondente a \(\eta(v) + 1\), qualsiasi questo sia. In un'altra \(\eta\) sarà un altro numero, che sarà \(\eta'(v) + 1\).

Quindi per essere precisi la definizione di \(\Eff\) è la seguente:
\begin{definition}{\(\Eff\) ben definita}
  \[\Eff(a,\eta)=\eta[v\leftarrow \Val(\expr, \eta)]\]
  con
  \[\eta[v\leftarrow \Val(\expr,\eta)] = \eta'\]
  dove
  \[\eta' (x)= \begin{cases}
      \Val(\expr,\eta) & \text{se } x = v  \\
      \eta(x)          & \text{altrimenti}
    \end{cases}\]
\end{definition}


\subsection{Relazione di Transizione}
Per definire ora la funzione di transizione dovremmo definire quali sono le coppie di stati che sono in relazione semantica, ovvero sono collegati da una transizione basandoci sulla combinazione di \(\hookrightarrow\) e \(\Eff\).

\begin{definition}{Relazione di transizione \(\xrightarrow{}\)}
  Definiamo la relazione di transizione \(\xrightarrow{} \subseteq S \times \Act \times S\) tale che:
  \begin{itemize}
    \item Se \(l \xhookrightarrow{g:a} l' \) è nel \(PG\) e \(\eta \models g\) allora \(\langle l,\eta \rangle \xrightarrow{a} \langle l',\eta' \rangle\) è nella relazione di transizione, con \(\eta' = \Eff(a,\eta)\).
  \end{itemize}
\end{definition}

\subsection{Stati iniziali}

\begin{definition}{\(S_0\)}
  Gli stati iniziali sono definiti come \(S_0 = \set{\langle l, \eta \rangle \mid l \in \Loc_0 \land \eta \models g_0}\)
\end{definition}

\subsection{Etichettatura degli stati}

\begin{definition}{\(AP\)}
  Per l'etichettatura, definiamo con \(AP = \Loc \cup \Cond(\Var)\) e \(L : S \to 2^{AP}\) tale che \(L(\langle l,\eta\rangle) = \set{l} \cup \set{g \mid g \in \Cond(\Var) \land {\eta \models g}}\).
\end{definition}

\section{Composizione e Sincronizzazione}

Fatto ciò, abbiamo definito il linguaggio per definire i componenti, ma questo è poco più di un programma sequenziale. Un po' di più perché da possibilità di scelte non deterministiche.
Per sistemi complessi è però molto difficile specificarli per intero.

Dunque l'approccio più comune è quello di specificarli tramite composizione di componenti più semplici, che interagiscono in modo sequenziale o concorrente.
La concorrenza è modellabile in tantissimi modi diversi.

\todo{Riccardo: Esempi da trascrizione (Dovrebbe stare nella parte 6 della trascrizione)}

\subsection{Composizione parallela sincrona}
La \keyword{composizione parallela sincrona} è facile da definire.

\begin{definition}{Composizione parallela sincrona}
  Dati \(TS_1, \ldots, TS_n\) sistemi a transizione tali che \(TS_i = (S_i, S_{0i}, A_i, R_i)\) con \(S_i\) disgiunti, allora la composizione parallela sincrona è il sistema a transizioni \(TS = (S,S_0,R,L)\) con:
  \begin{itemize}
    \item \(S = S_1 \times \cdots \times S_n\), stati globali
    \item \(S_0 = S_{01} \times \cdots \times S_{0n}\),
    \item \(A = A_1 \cup \ldots \cup A_n\), azioni globali
    \item \(R \subseteq S \times A \times S\) è tale che, dati \(\langle s_1, \ldots, s_n\rangle, \langle s_1', \ldots, s_n'\rangle \in S\) e \(\langle a_1,\ldots, a_n \rangle \in A\) contiene \(\langle s_1,\ldots,s_n \rangle \xrightarrow{\langle a_1,\ldots,a_n \rangle} \langle s_1',\ldots,s_n'\rangle\) se \(s_i \xrightarrow{a_i} s_i'\,, \forall i \in \set{1,\ldots,n}\).
  \end{itemize}
\end{definition}

Prendiamo come esempio una coppia di Transition Systems, che rappresentano due astrazioni di semafori.

Il semaforo si può trovare in due stati, rosso o verde. Inoltre ci sono due transizioni: una che passano dallo stato rosso allo stato verde etichettata con \q{OK} e l'altra che passa dallo stato verde allo stato rosso etichettata con \q{KO}.

Immaginando di avere due semafori con stato iniziale differente otteniamo la seguente figura.

\begin{figure}[H]
  \centering
  \begin{tikzpicture}[shorten >=1pt,node distance=3cm,on grid,auto, bend right=30]
    \node[state,initial, label=\(TS_1\)]  (RED)                      {RED};
    \node[state]          (GREEN) [below=of RED] {GREEN};

    \path[->] (RED) edge[bend right] node[swap] {OK} (GREEN)
    (GREEN) edge[bend right] node[swap] {KO} (RED);

    \node[state, initial, label=\(TS_2\)] (GREEN) [right=7cm of RED] {GREEN};
    \node[state]  (RED)   [below=of GREEN]  {RED};

    \path[->] (RED) edge[bend right] node[swap] {OK} (GREEN)
    (GREEN) edge[bend right] node[swap] {KO} (RED);
  \end{tikzpicture}
  \caption{Transition Systems di due semafori}\label{fig:ts_semaphore}
\end{figure}

Se adesso applichiamo l'operazione di composizione sincrona ai due Transition Systems otteniamo il seguente:

\begin{figure}[H]
  \centering
  \begin{tikzpicture}[shorten >=1pt,node distance=4cm,on grid,auto, bend right=30]
    \node[state,initial]  (RED-GREEN) {RED,GREEN};
    \node[state] (GREEN-RED) [below=of RED-GREEN] {GREEN,RED};

    \path[->] (RED-GREEN) edge[bend right] node[swap] {\(\langle OK,KO\rangle\)} (GREEN-RED)
    (GREEN-RED) edge[bend right] node[swap] {\(\langle KO,OK\rangle\)} (RED-GREEN);

    \node[draw=none, font=\Huge] (label) [above=2cm of RED-GREEN] {\(TS_1\,\|\,TS_2\)};

  \end{tikzpicture}
  \caption{Composizione sincrona di Transition Systems di due semafori}\label{fig:ts_semaphore_sync_composition}
\end{figure}

In realtà il risultato della composizione non è proprio così, poiché gli stati derivanti dal prodotto cartesiano sono quattro: (red, green), (red, red), (green, green), (green red) come riportato dalla figura è la seguente:

\begin{figure}[H]
  \centering
  \begin{tikzpicture}[shorten >=1pt,node distance=4cm,on grid,auto, bend right=30]
    \node[state,initial]  (RED-GREEN) {RED,GREEN};
    \node[state] (GREEN-RED) [below=of RED-GREEN] {GREEN,RED};

    \path[->] (RED-GREEN) edge[bend right] node[swap] {\(\langle OK,KO\rangle\)} (GREEN-RED)
    (GREEN-RED) edge[bend right] node[swap] {\(\langle KO,OK\rangle\)} (RED-GREEN);

    \node[state] (RED-RED) [right=7cm of RED-GREEN] {RED,RED};
    \node[state] (GREEN-GREEN) [below=of RED-RED] {GREEN,GREEN};

    \path[->] (RED-RED) edge[bend right] node[swap] {\(\langle OK,OK\rangle\)} (GREEN-GREEN)
    (GREEN-GREEN) edge[bend right] node[swap] {\(\langle KO,KO\rangle\)} (RED-RED);

    \node[draw=none, font=\Huge] (label) [above=2cm of RED-GREEN, xshift=4cm] {\(TS_1\,\|\,TS_2\)};
  \end{tikzpicture}
  \caption{Composizione sincrona di Transition Systems di due semafori (completa)}\label{fig:ts_semaphore_sync_composition_complete}
\end{figure}

\begin{note}{}
  Notiamo però che non essendoci percorsi dallo stato iniziale (RED,GREEN) agli altri due stati, li possiamo escludere poiché stati che non saranno mai raggiungibili.
\end{note}

\subsection{Composizione parallela asincrona}

Se da un lato ci sono i modelli completamenti sincroni dall'altra ci sono quelli completamente asincroni, che generalizzano quelli in cui alcuni si sincronizzano e altri no.

L'idea è quella di avere un modello che permette di descrivere sistemi che evolvono indipendentemente da eventi esterni.

Per farlo in maniera semplice possiamo usare il modello di composizione sincrona ma aggiungendo ad ogni insieme di azioni locali una \textit{no operation} \(-\) per permettere a ogni automa locale di effettuare un'azione vacua, ovvero un azione che non modifica lo stato del sistema.


In generale la composizione sincrona è definita come la composizione sincrona, ma con l'aggiunta dell'azione \textit{no operation}, in cui intuitivamente se ci troviamo in un particolare stato e abbiamo come  transizione da eseguire l'azione \q{\(-\)} allora lo stato partenza sarà uguale allo stato di arrivo.

Questo permette, inoltre, a stati che precedentemente non erano raggiungibili di diventarlo, poiché non vi è più una sincronizzazione esterna.

Definiamo adesso formalmente la composizione asincrona.

\begin{definition}{Composizione asincrona}
  Assumiamo \(TS_1,\ldots,TS_n\) Transition Systems tali che \(\forall i \in 1\ldots n, TS_i = \langle S_i,A_i,R_i,S_{i0}\rangle \). Il Transition system ottenuto dalla composizione asincrona di questi transition systems è \( TS = TS_1 \| \ldots \| TS_n = \langle S,A,R,S_0\rangle \) tale che:
  \begin{itemize}
    \item \( S = S_1 \times \cdots \times S_n \)
    \item \( A = (A_1 \cup \{-\}) \times \cdots \times (A_n \cup \{-\}) \)
    \item \( S_0 = S_{10} \times \cdots \times S_{n0} \)
    \item \( R \) contiene \( \langle s_1,\ldots,s_n \rangle \xrightarrow{\langle a_1,\ldots,a_n \rangle} \langle s_1',\ldots,s_n' \rangle \), se \( \langle a_1,\ldots,a_n \rangle \in A \) e, per ogni \( 1 \leq i \leq n \),
          \begin{itemize}
            \item \( s_i = s_i' \) e \( a_i = - \) oppure
            \item \( s_i \xrightarrow{a_i}_i s_i' \) e \( a_i \neq - \).
          \end{itemize}
  \end{itemize}
\end{definition}

Questo modello ammette vari tipi di comportamenti, da comportamenti sincroni a comportamenti del tutto asincroni,

Inoltre questo tipo di modello in cui la concorrenza è virtuale, è esattamente come vengono gestiti nella pratica i processi concorrenti, infatti come sott'insieme di questo modello abbiamo la composizione asincrona con interleaving, in cui c'è composizione completamente asincrona, ovvero in ogni momento un solo componente esegue un'azione, ma tale scelta è non deterministica.

Nel caso dello schedule dei processi, la scelta non è non deterministica, infatti ci sono diverse politiche che possono essere usate come criterio di scelta. Ma in generale, nel caso in cui non si riuscisse a definire una politica da adottare, la scelta sarà non deterministica, come infatti avviene nei Labeled Transition System.

Dunque mentre la composizione sincrona non aggiunge nessun non determinismo se non già presente negli LTS che lo compongono, la composizione asincrona interleaving e quella generale introducono implicitamente scelte non determistiche.

Intuitivamente come definizione possiamo usare quella generale, a cui applichiamo la seguente restrizione: ogni azione globale è composta da una sola azione locale e no operation per gli altri stati in maniera non deterministica.


\begin{definition}{Composizione Asincrona Interleaving}
  Assumiamo \(TS_1,\ldots,TS_n\) Transition Systems tali che \(\forall i \in 1\ldots n, TS_i = \langle S_i,A_i,R_i,S_{i0}\rangle \). Il Transition system ottenuto dalla composizione asincrona (interleaving) di questi transition systems è \( TS = TS_1 \parallel \ldots \parallel TS_n = \langle S,A,R,S_0\rangle \) tale che:
  \begin{itemize}
    \item \( S = S_1 \times \cdots \times S_n \)
    \item \( A \subset (A_1 \cup \{-\}) \times \cdots \times (A_n \cup \{-\}) \) tale che \( \langle a_1,\ldots,a_n \rangle \in A \) se e solo se, per ogni \( j \neq i \), \( a_i \in A_i \) implica \( a_j = - \)
    \item \( S_0 = S_{10} \times \cdots \times S_{n0} \)
    \item \( R \) contiene \( \langle s_1,\ldots,s_n \rangle \xrightarrow{\langle a_1,\ldots,a_n \rangle} \langle s_1',\ldots,s_n' \rangle \), se \( \langle a_1,\ldots,a_n \rangle \in A \) e, per ogni \( 1 \leq i \leq n \),
          \begin{itemize}
            \item \( s_i = s_i' \) e \( a_i = - \) oppure
            \item \( s_i \xrightarrow{a_i}_i s_i' \) e \( a_i \neq - \).
          \end{itemize}
  \end{itemize}
\end{definition}

Ovviamente in questo modello non è inficiata la raggiungibilità di stati, ma non vi è vero parallelismo, ed inoltre vi è una perdita delle informazioni riguardo al numero dei percorsi e quindi di computazioni. In generale però è un modello che viene utilizzato molto frequentemente poiché offre una buona astrazione.

Consideriamo l'esempio precendente e applichiamo la composizione asincrona interleaving:

\begin{figure}[H]
  \centering
  \begin{tikzpicture}[shorten >=1pt,node distance=3cm,on grid,auto, bend right=30]
    \node[state,initial] (RED-GREEN) {RED,GREEN};
    \node[state] (GREEN-RED) [right=7cm of RED-GREEN] {GREEN,RED};
    \node[state] (RED-RED) [above right=5cm of RED-GREEN] {RED,RED};
    \node[state] (GREEN-GREEN) [below right=5cm of RED-GREEN] {GREEN,GREEN};
    \node[draw=none, above=1.5cm of RED-RED] {};


    \path[->] (RED-GREEN) edge[bend right] node[] {\(\langle -,KO\rangle\)} (RED-RED)
    edge[bend right] node[swap] {\(\langle OK,-\rangle\)} (GREEN-GREEN)
    (RED-RED) edge[bend right] node[swap] {\(\langle -,OK\rangle\)} (RED-GREEN)
    (GREEN-GREEN) edge[bend right] node[] {\(\langle KO,-\rangle\)} (RED-GREEN)
    ;


    \path[->] (GREEN-RED) edge[bend right] node[swap] {\(\langle KO,-\rangle\)} (RED-RED)
    edge[bend right] node[] {\(\langle -,KO\rangle\)} (GREEN-GREEN)
    (RED-RED) edge[bend right] node[] {\(\langle OK,-\rangle\)} (GREEN-RED)
    (GREEN-GREEN) edge[bend right] node[swap] {\(\langle -,KO\rangle\)} (GREEN-RED)
    ;
  \end{tikzpicture}
  \caption{Composizione asincrona interleaving di Transition Systems di due semafori}\label{fig:ts_semaphore_async_composition}
\end{figure}


Una formulazione equivalente del modello asincrono con interleaving più compatto è quella di avere come insieme delle azioni l'unione disgiunta degli insiemi di azioni locali, ad esempio avendo le azioni rappresentate come coppie \(\langle i,a\rangle\) con \(i\) indice dell'automa locale e \(a\) azione locale. Allora possiamo appropriamente riformulare la relazione di transizione nel seguente modo.

\begin{definition}{Composizione Asincrona Interleaving Compatta}
  Assumiamo \(TS_1,\ldots,TS_n\) Transition Systems tali che \(\forall i \in 1\ldots n, TS_i = \langle S_i,A_i,R_i,S_{i0}\rangle \). Il Transition system ottenuto dalla composizione asincrona (interleaving) di questi transition systems è \( TS = TS_1 \parallel \ldots \parallel TS_n = \langle S,A,R,S_0\rangle \) tale che:
  \begin{itemize}
    \item \( S = S_1 \times \cdots \times S_n \)
    \item \( A = \bigcup_{1 \leq i \leq n} A_i \) (oppure l'unione disgiunta \(\uplus\))
    \item \( S_0 = S_{10} \times \cdots \times S_{n0} \)
    \item \( R \) contiene \( \langle s_1,\ldots,s_n \rangle \xrightarrow{a} \langle s_1,\ldots,s_i',\ldots,s_n \rangle \), se si ha \( s_i \xrightarrow{a}_i s_i' \), per qualche \( 1 \leq i \leq n \).
  \end{itemize}
\end{definition}


Noi avevamo però dei Program Graph per descrivere i componenti, che poi venivano \q{compilati} in LTS\@. Abbiamo visto come fare composizione di LTS, ma per i program graph dobbiamo modellare la possibilità di poter condividere informazioni, modellare forme di comunicazioni tra processi. Uno dei modi per fare ciò è attraverso la condivisione di variabili, ovvero condividere aree di memoria in cui più processi possono leggere e scrivere.

Nel mondo degli LTS però non c'è il concetto di variabile. Quindi l'approccio di tradurre vari PG in LTS separatamente per poi comporli non è necessariamente corretto, poiché nel passaggio da PG a LTS si perde il concetto di variabile e dunque di condivisione di variaibli.

\begin{example}{}
  Vediamo cosa succede nello specifico con un esempio.
  \begin{figure}[H]
    \centering
    \includegraphics[width=\textwidth]{_images/pg_incorrect_variable_sharing_example}
    \caption{Esempio di condivisione delle variabili errata}\label{fig:pg_incorrect_variable_sharing_example}
  \end{figure}
\end{example}

Quindi quello che succede è che quando ottengo l'LTS composto, non vi è nessuna releazione tra le due variabili \(x\) quindi operazioni come incremento o decremento, esse avvengono soltanto sull'oggetto su cui si applicano.

Dunque questa fusione va fatta a livello di PG, ma allora la composizione va fatta a livello di PG, e dunque è necessario definire una composizione di PG\@.

\subsection{Interleaving Program Graph}

\begin{definition}{Composizione Asincrona (Interleaving) di Program Graph}
  Dati \(n\) Program Graph \(PG_i = (\Loc_i, \Act_i, \Eff_i, \hookrightarrow_i, \Loc_{i0}, g_{i0})\) definiti sui rispettivi insiemi di variabili \(\Var_i\), il Program Graph ottenuto dalla loro composizione asincrona con interleaving è:
  \[ PG = PG_1 \parallel PG_2 \parallel \ldots \parallel PG_n = (\Loc, \Act, \Eff, \hookrightarrow, \Loc_0, g_0) \]
  definito sull'insieme di variabili globali \(\Var = \bigcup_{1 \leq i \leq n} \Var_i\), dove:
  \begin{itemize}
    \item \(\Loc = \Loc_1 \times \Loc_2 \times \cdots \times \Loc_n\) è il prodotto cartesiano delle locazioni;
    \item \(\Act = \bigcup{1 \leq i \leq n} \Act_i\) è l'unione delle azioni (può essere anche disgiunta);
    \item \(\Loc_0 = \Loc_{10} \times \Loc_{20} \times \cdots \times \Loc_{n0}\) è l'insieme delle locazioni iniziali;
    \item \(g_0 = g_{10} \land g_{20} \land \ldots \land g_{n0}\) è la condizione (o guardia) iniziale globale.
  \end{itemize}
\end{definition}

\paragraph{Condivisione delle variabili:}
L'insieme delle variabili globali \(\Var\) è l'unione (non disgiunta) delle variabili dei singoli processi. Per semplicità, assumiamo che se due variabili locali in processi diversi sono definite con lo stesso nome, esse siano implicitamente condivise\footnote{Nel caso in cui due variabili omonime non debbano essere condivise, è necessario disambiguarle prima della composizione (ad esempio rinominandole). Una definizione alternativa potrebbe tenere conto di variabili non condivise con lo stesso nome, ma risulterebbe inutilmente più complessa senza aggiungere effettivo potere espressivo.}.
In questo modo, il problema riscontrato con l'LTS globale (in cui si perdeva la relazione tra le variabili) non sussiste più: la modifica che un processo opera su una variabile condivisa diventa immediatamente visibile a tutti gli altri processi.

\paragraph{Relazione di Transizione (\(\hookrightarrow\)):}
Le transizioni vengono definite in modo del tutto analogo all'interleaving degli LTS\@. Se il processo \(i\)-esimo esegue un'azione locale, la locazione globale si aggiorna modificando solo la componente \(i\)-esima, lasciando inalterate le locazioni di tutti gli altri processi. Formalmente:
\[
  \text{Se } l_i \xhookrightarrow{g:a}_i l_i' \text{ in } PG_i \implies \langle l_1, \ldots, l_i, \ldots, l_n \rangle \xhookrightarrow{g:a} \langle l_1, \ldots, l_i', \ldots, l_n \rangle
\]

\paragraph{Definizione della funzione degli effetti (\(\Eff\)) globale:}
L'aspetto più complesso della composizione riguarda la funzione \(\Eff\), che funge da ponte tra la sintassi e la semantica. Dobbiamo infatti costruire una funzione globale
\[ \Eff : \Act \times \Eval(\Var) \to \Eval(\Var) \]
avendo a disposizione esclusivamente le funzioni locali
\[ \Eff_i : \Act_i \times \Eval(\Var_i) \to \Eval(\Var_i) \]

Presa un'azione arbitraria \(a \in \Act\) e una valutazione globale \(\eta \in \Eval(\Var)\), dobbiamo costruire la nuova valutazione globale \(\eta'\). Essendo \(\Act\) un'unione disgiunta, sappiamo esattamente a quale set locale \(\Act_i\) appartiene l'azione \(a\).

Tuttavia, non possiamo passare direttamente \(\eta\) alla funzione \(\Eff_i\), poiché i domini sono disallineati: \(\eta\) è valutata su tutto \(\Var\), mentre \(\Eff_i\) richiede una valutazione definita solo su \(\Var_i\).
Per risolvere questo problema, procediamo in tre fasi:
\begin{enumerate}
  \item \textbf{Restrizione:} Definiamo una valutazione locale \(\eta_i\) ottenuta restringendo il dominio della funzione globale \(\eta\) alle sole variabili del processo \(i\)-esimo. Ovvero: \(\eta_i(v) = \eta(v)\) per ogni \(v \in \Var_i\).
  \item \textbf{Applicazione:} Ora che i formati sono compatibili, possiamo applicare la funzione locale per calcolare i nuovi valori delle variabili modificate dall'azione: \(\eta_i' = \Eff_i(a, \eta_i)\).
  \item \textbf{Estensione (Completamento):} Infine, ricostruiamo la nuova valutazione globale \(\eta'\). Per le variabili appartenenti a \(\Var_i\), assegniamo il risultato appena calcolato \(\eta_i'\). Per tutte le altre variabili di sistema che non appartengono a \(\Var_i\), applichiamo sostanzialmente la funzione identità, mantenendo il valore che avevano nella \(\eta\) originale.
\end{enumerate}

Formalmente, l'effetto di un'azione sulla valutazione \(\eta\) per una qualsiasi variabile \(v \in \Var\) si definisce come:
\[
  \Eff(a, \eta)(v) =
  \begin{cases}
    \Eff_i(a, \eta_i)(v) & \text{se } a \in \Act_i \text{ e } v \in \Var_i \\
    \eta(v)              & \text{altrimenti}
  \end{cases}
\]
dove \(\eta_i\) è la restrizione di \(\eta\) alle variabili in \(\Var_i\).


\begin{example}{Composizione di PG con Variabile Condivisa}
  Riprendendo il problema della variabile condivisa visto in \Cref{fig:pg_incorrect_variable_sharing_example}, se effettuiamo la composizione a livello di Program Graph otteniamo un modello corretto \(PG_1 \parallel PG_2\).

  Nello stato iniziale globale \((l_1^1, l_1^2)\), una singola valutazione \(\eta\) traccia il valore dell'unica variabile condivisa \(x\).
  Se il primo processo esegue l'incremento (\(x := x+1\)), il sistema transita nella locazione globale \((l_2^1, l_1^2)\). Successivamente, se il secondo processo esegue il decremento (\(x := x-1\)), il sistema raggiunge lo stato \((l_2^1, l_2^2)\). Applicando la semantica della funzione \(\Eff\) globale, queste operazioni avvengono sulla medesima istanza di \(x\), riportandola coerentemente al suo valore originale. Lo stesso vale per l'ordine inverso di esecuzione.
\end{example}

Dal momento che non c'è nessun impegno semantico a modellare cosa succede veramente a livello puramente sintattico, il comportamento reale emerge durante la traduzione in LTS (utilizzando la funzione \(\Eff\) globale come ponte).
Durante questa compilazione, gli stati dell'LTS diventano coppie formate dalla locazione globale e dalla valutazione globale delle variabili: \(\langle (l_a, l_b), \eta \rangle\).

In assenza di variabili condivise, la composizione diretta degli LTS originali non causerebbe problemi poiché i domini sarebbero del tutto indipendenti. Il problema sorge unicamente in presenza di memoria condivisa: se gli LTS non mantengono traccia di questa condivisione, le modifiche rimangono isolate alla realtà di un singolo processo, rendendo impossibile ristabilire una coerenza nello stato comune. L'approccio corretto è dunque spostare l'operazione di composizione parallelamente in alto, a livello del linguaggio astratto (Program Graph).


\subsection{Protocolli di Mutua Esclusione}

Per calare questi concetti di modellazione in un contesto pratico, introdurremo degli esempi concreti ispirati ai \keyword{protocolli di mutua esclusione}.

Questi sistemi sono sufficientemente semplici ma ricchi di sfumature da prestarsi perfettamente allo studio con i tool di verifica formale (Model Checking). Tali strumenti, infatti, storicamente nascono proprio per analizzare ed evidenziare difetti in questa tipologia di protocolli.

I protocolli di mutua esclusione che analizzeremo sono progettati per operare in modo \keyword{distribuito}, ovvero senza un sistema centralizzato di risoluzione dei conflitti. La sincronizzazione avviene esclusivamente ed in maniera autonoma tra i processi tramite \keyword{variabili condivise}.

\begin{example}{Un Primo Tentativo di Protocollo}
  \begin{figure}[H]
    \centering
    \includegraphics[width=\textwidth]{_images/mutual_exclusion_protocol_code_example}
    \caption{Esempio di Protocollo di Mutua Esclusione}\label{fig:mutual_exclusion_protocol_example}
  \end{figure}

  Questo primo tentativo garantisce la mutua esclusione tramite l'uso di due variabili booleane condivise, \texttt{flag1} e \texttt{flag2}. Ogni processo utilizza il proprio flag per segnalare l'intenzione di accedere alla sezione critica.
\end{example}

Nota che, l'obiettivo della verifica formale in questo contesto non è analizzare lo specifico codice di un processo applicativo, bensì il \keyword{meccanismo di mutua esclusione} in sé. Per questo motivo, si ricorre a una forte astrazione: la logica del protocollo di comunicazione rimane esplicita nel modello, mentre il resto del codice (che varia da processo a processo ed è arbitrariamente lungo) viene astratto in blocchi generici e opachi.

Strutturalmente, ogni processo esegue ciclicamente le seguenti fasi:
\begin{enumerate}
  \item Sezione di codice non critico (\textit{non-critical code});
  \item Protocollo di ingresso (\textit{entry protocol});
  \item Sezione critica (\textit{critical section});
  \item Protocollo di uscita (\textit{exit protocol});
\end{enumerate}


Per poter sottoporre questo codice a verifica formale, dobbiamo prima modellarlo. Possiamo astrarre le diverse fasi dell'esecuzione assegnando delle locazioni specifiche a ciascuna sezione di codice:
\begin{itemize}
  \item \(nc_a, nc_b\): astrae e racchiude tutta la parte di esecuzione del codice non critico indipendente.
  \item \(wt_a, wt_b\): attesa per l'accesso, o fase di \textit{entry protocol} (in cui si testa il while);
  \item \(cr_a, cr_b\): esecuzione della sezione critica (\textit{critical section});
  \item \(ex_a, ex_b\): rilascio della risorsa, o fase di \textit{exit protocol}.
\end{itemize}

Dunque un possibile automata per il processo A e B è il seguente:

\begin{figure}[H]
  \centering
  \includegraphics[width=\textwidth]{_images/automata_for_mutual_exclusion}
  \caption{Automata del Processo A e B per l'esempio di protocollo di Mutua Esclusione}\label{fig:automata_mutual_exclusion_protocol_example}
\end{figure}

Andiamo adesso ad analizzare le transizioni per il processo A

\begin{itemize}
  \item \(\mathbf{nc_a \xrightarrow{\text{True: } flag_1 := \text{True}} wt_a}\): Dalla sezione non critica, il processo può decidere in qualsiasi momento (guardia \texttt{True}) di avviare il protocollo di ingresso, impostando il proprio flag a \texttt{True} e portandosi nello stato di attesa.
  \item \(\mathbf{wt_a \xrightarrow{flag_2} wt_a}\): Se il flag del processo avversario è \texttt{True}, il processo rimane bloccato nello stato d'attesa (attesa attiva).
  \item \(\mathbf{wt_a \xrightarrow{\neg flag_2} cr_a}\): Se il flag avversario è \texttt{False}, la sezione critica è libera e il processo vi accede.\ \textit{Nota:} questa guardia e la precedente sono ovviamente mutuamente esclusive.
  \item \(\mathbf{cr_a \xrightarrow{\text{True}} ex_a}\): L'uscita dalla sezione critica può avvenire in qualsiasi momento, portando il sistema nella fase di exit.
  \item \(\mathbf{ex_a \xrightarrow{\text{True: } flag_1 := \text{False}} nc_a}\): Il processo rilascia la risorsa impostando il proprio flag a \texttt{False} e ritorna in esecuzione non critica.
\end{itemize}

\begin{note}{Assunzione di Atomicità}
  Questa versione astratta assume implicitamente che l'intera esecuzione della porzione di codice associata ad una locazione (\(nc, wt,\ldots\)) si esaurisca in un singolo passo atomico del modello.
\end{note}

Andando adesso a comporre i due automata, si ottiene l'LTS globale del sistema. Negli stati di questo LTS, la funzione di etichettatura traccerà le locazioni correnti di entrambi i processi e le variabili booleane vere in quel momento (se una proposizione non compare, è falsa).

Ad esempio, lo stato iniziale è denotato come \((nc_a, nc_b)\), senza flag, il che significa che entrambi i processi sono nel codice non critico e \(flag_1, flag_2\) sono entrambi falsi.

Ecco adesso un frammento dell'LTS finale:

\begin{figure}[H]
  \centering
  \includegraphics[width=\textwidth]{_images/lts_fragment_for_mutual_exclusion}
  \caption{Frammento dell'LTS per l'esempio di protocollo di Mutua Esclusione}\label{fig:lts_fragment_mutual_exclusion_protocol_example}
\end{figure}


Analizzando le possibili computazioni in questo LTS, emergono immediatamente due criticità fondamentali:

\begin{itemize}
  \item Situazione di Deadlock:
        Esiste una traccia in cui il Processo A esegue il primo passo portandosi in \(wt_a\) (e impostando \(flag_1\) a \texttt{True}). Subito dopo, lo schedulatore cede il controllo al Processo B, il quale a sua volta esegue il primo passo portandosi in \(wt_b\) (impostando \(flag_2\) a \texttt{True}). (Nel LTS finale non viene rappresentato questo stato poiché è un frammento, ma nell'LTS completo esso compare)
        A questo punto, il sistema si trova nello stato globale \((wt_a, wt_b)\) con \(f_1\) e \(f_2\) entrambi veri. Nessuno dei due processi potrà mai procedere verso la sezione critica: il sistema è in \keyword{deadlock}. Dunque il protocollo descritto è errato.
  \item Starvation (Esecuzioni Spurie):
        Esiste inoltre una traccia circolare infinita in cui il Processo A viene ininterrottamente scelto per l'esecuzione:
        \[
          (nc_a, nc_b) \xrightarrow{f_1:=T} (wt_a, nc_b) \xrightarrow{\neg f_2} (cr_a, nc_b) \xrightarrow{} (ex_a, nc_b) \xrightarrow{f_1:=F} (nc_a, nc_b) \dots
        \]
        In questa sequenza, il Processo A entra nella sezione critica, esce, e vi rientra ripetutamente all'infinito, mentre il Processo B rimane bloccato in \(nc_b\) senza mai avere la possibilità di agire.

        Perché si verifica questo comportamento indesiderato? Il problema non è intrinseco al protocollo di mutua esclusione, ma è una conseguenza diretta di come stiamo modellando l'interazione concorrente.
\end{itemize}


\subsection{Lo Schedulatore e il Concetto di Fairness}

Nel nostro modello LTS, lo schedulatore dei processi è modellato implicitamente tramite il non-determinismo puro.

Il non-determinismo puro assume che, se ci sono più transizioni possibili, qualsiasi di esse può essere scelta, senza alcun vincolo. Ciò significa che il modello ammette uno schedulatore \q{non fair} che sceglie sistematicamente e all'infinito di eseguire il Processo A, ignorando il Processo B.

Per risolvere questa anomalia potremmo essere tentati di modellare esplicitamente uno specifico schedulatore aggiungendolo come componente al sistema. Tuttavia, questo approccio va contro l'idea di avere un modello astratto, poiché quello che si vorrebbe verificare è che il protocollo di mutua esclusione sia corretto \textit{a prescindere} dalla specifica politica di schedulazione utilizzata, a patto che questa sia ragionevole.

Per risolvere questo problema senza appesantire il modello con la specifica dettagliata di uno scheduler, si introduce il concetto di \keyword{fairness} (equità). A livello di verifica formale, la fairness non viene modellata modificando l'LTS, ma agisce come un filtro sulle esecuzioni.

L'intuizione è la seguente: l'algoritmo di analisi non deve considerare tutti i percorsi possibili dell'LTS, ma solo quelli definiti \textit{fair}. I percorsi che violano i vincoli di equità vengono ignorati, poiché rappresentano comportamenti che, sebbene strutturalmente possibili nel grafo, non sono concretamente realizzabili in un sistema operativo reale. Questo ci permette di verificare le proprietà del protocollo sotto qualsiasi scheduler \q{ragionevole}

\begin{definition}{I Vincoli di Fairness}
  Esistono tre categorie principali di vincoli di fairness, con diversi gradi di restrizione:

  \begin{description}
    \item[Unconditional Fairness:] Ogni processo deve essere schedulato infinitamente spesso. In altre parole, non può esistere un punto nell'esecuzione oltre il quale un processo non viene mai più scelto.
      \begin{itemize}
        \item \textit{Limite:} Nel contesto degli scheduler, è un vincolo troppo debole. Non garantisce che il processo evolva, poiché potrebbe essere schedulato proprio quando tutte le sue azioni sono disabilitate dalle guardie.
      \end{itemize}

    \item[Weak Fairness:] Se un processo è \textit{continuamente} abilitato, allora deve essere schedulato infinitamente spesso.
      \begin{itemize}
        \item \textit{Limite:} È ancora piuttosto restrittiva, poiché richiede che l'abilitazione non venga mai meno affinché il vincolo si applichi.
      \end{itemize}

    \item[Strong Fairness:] Se un processo è abilitato \textit{infinitamente spesso} (anche se non continuativamente), allora deve essere schedulato infinitamente spesso.
      \begin{itemize}
        \item \textit{Vantaggio:} È il vincolo più ragionevole per la modellazione di sistemi concorrenti, poiché copre anche i casi in cui un processo alterna momenti di abilitazione a momenti di blocco.
      \end{itemize}
  \end{description}
\end{definition}


Nel caso del protocollo di mutua esclusione analizzato precedentemente, uno qualsiasi di questi tre vincoli sarebbe sufficiente per escludere il ciclo infinito in cui solo il Processo A accede alla risorsa. Infatti, mentre A cicla, il Processo B si trova in una locazione (come \(nc_b\)) dove ha sempre almeno un'azione abilitata (avendo guardia \texttt{true}).

% NEW_LESSON
Resta però il problema di Deadlock per questo prtocollo. Poiché, se il sistema esegue un interleaving tale per cui entrambi i processi impostano il proprio flag a \code{true} prima che uno dei due riesca a valutare la condizione di ingresso, si genera uno stato dal quale nessuno dei due può più procedere. Entrambi rimarranno bloccati nei rispettivi cicli di attesa.

\begin{figure}{H}
  \centering
  \includegraphics[width=\textwidth]{_images/lts_fragment_deadlock_for_mutual_exclusion}
  \caption{Deadlock nell'LTS dell'esempio di protocollo di Mutua Esclusione}\label{fig:lts_deadlock_mutual_exclusion_protocol_example}
\end{figure}


Per superare i limiti e le vulnerabilità al \keyword{deadlock} di un protocollo rudimentale basato solo sui flag, si possono impiegare soluzioni più sofisticate, come l'algoritmo di Peterson.


\subsection{Astrazione Temporale e Generalizzazione del Modello}

Tralasciando però questo problema che si può dunque risolvere adottando un algoritmo più furbo, un'altro problema è che il modello costruito è piuttosto limitativo poiché stiamo modellando processi in cui il codice critico e non critico è compattato in un'unica istruzione, e dunque il tempo che intercorre tra le varie sezioni è unitaria.

Questo infatti potrebbe portare ad una violazione di qualche proprietà se il processo resta più di un instante di tempo in una particolare sezione.

Tuttavia, l'obiettivo della verifica formale non è modellare l'esatta esecuzione temporale di un'istanza software, bensì validare la robustezza logica del protocollo per \ti{qualsiasi} suo potenziale utilizzatore. L'unica assunzione rigorosa e necessaria affinché proprietà come l'assenza di starvation siano formalmente valide è che ogni processo permanga per un tempo \keyword{finito} nelle proprie sezioni di codice. Se si ammettesse l'eventualità che un processo acquisisca la risorsa per poi trattenersi in un ciclo infinito nella sezione critica, l'insuccesso dell'altro processo dipenderebbe da un difetto semantico del programma utente (che non rilascia la risorsa) e non da una falla del protocollo di sincronizzazione. Sotto la garanzia di permanenza finita, le proprietà verificate sul modello astratto si mantengono inalterate indipendentemente dalla durata reale dei task.


Per verificare la correttezza di un protocollo per una classe infinita di potenziali processi, è necessario costruire un modello generalizzato che astragga dai dettagli implementativi del codice utente. In un sistema a transizioni a tempo discreto, l'unità elementare di tempo coincide con l'esecuzione di una singola transizione, che quindi comporta che le azioni siano atomiche, ovvero che l'esecuzione di un'azione non sia scomponibile ulteriormente.

Per modellare il fatto che un processo possa trascorrere una quantità arbitraria di tempo all'interno di una specifica sezione di codice (critica o non critica), introduciamo esplicitamente il \keyword{non determinismo} aggiungendo delle transizioni \q{ciclo} (self-loop) sugli stati di nostro interesse. Ad esempio, possiamo aggiungere transizioni etichettate incondizionatamente con \code{True} sugli stati \(nc_a\) (codice non critico) e \(cr_a\) (sezione critica), permettendo all'automa di permanere nello stesso stato pur facendo evolvere il sistema globale.

Questo è importante soprattutto quando consideriamo la composizione tra più processi, poiché permette ad un altro processo di evolvere anche se il processo che stiamo considerando non modifica il proprio stato.

\subsection{Computazioni Spurie e Generalizzazione della Fairness}

L'introduzione dei self-loop risolve il problema della generalizzazione temporale ma espande eccessivamente i comportamenti possibili del modello, introducendo il rischio di \keyword{computazioni spurie}, come per esempio un ciclo infinito nella sezione critica di un processo.

\begin{definition}{Computazione Spuria}
  Una computazione è detta \ti{spuria} se è ammessa dal modello, ma non è realizzabile o ammissibile nel sistema concreto reale.
\end{definition}

Per eliminare le computazioni spurie e ripristinare il corretto funzionamento del modello senza rimuovere il non determinismo, si ricorre ancora una volta ai vincoli di \keyword{fairness}. I concetti di fairness discussi in precedenza per lo scheduling (unconditional, weak, strong) possono essere generalizzati ad ammettere \ti{qualsiasi proprietà di stato}, non solo l'abilitazione di un'azione.

Per mitigare l'anomalia del loop in sezione critica, è sufficiente imporre un vincolo di \keyword{unconditional fairness} basato sugli stati ovvero, che ogni processo debba rispettare il vincolo di non trovarsi in sezione critica per un tempo infinito.

Questo vincolo non impone alcun tetto massimo prefissato al numero di esecuzioni del self-loop (mantenendo quindi l'astrazione e permettendo un tempo arbitrario), ma impone rigorosamente che tale tempo sia finito.

\begin{note}{Applicazioni del non determinismo}
  Il non determinismo può essere utilizzato per modellare sistemi con informazioni parziali ovvero, permette di rappresentare decisioni interne del sistema che non sono note a priori, o comportamenti imprevedibili dell'ambiente esterno.
\end{note}

\subsection{Atomicità e Race Conditions}

Un aspetto critico durante la fase di modellazione è l'assunzione relativa all'\keyword{atomicità} delle operazioni. Nei Transition System (e nei Program Graph), una transizione rappresenta l'unità elementare di tempo e di evoluzione del sistema, incapsulando per definizione un'azione atomica (cioè un'operazione non ulteriormente scomponibile o interrompibile).

È fondamentale garantire che il livello di astrazione del modello rifletta accuratamente l'effettiva granularità del sistema concreto che si sta analizzando. In caso contrario, l'analisi potrebbe ignorare l'insorgenza di \keyword{race conditions} (condizioni di corsa), ovvero anomalie in cui l'esito della computazione dipende dall'ordine in cui i processi paralleli vengono schedulati.

\begin{example}{Incremento Concorrente e Race Condition}
  Si consideri un sistema con una variabile condivisa \(a\), inizializzata a \(0\), su cui operano contemporaneamente due processi eseguendo la medesima istruzione:
  \begin{center}
    Processo 1: \code{a = a + 1} \qquad Processo 2: \code{a = a + 1}
  \end{center}
  Se l'operazione \code{a = a + 1} è supportata a livello hardware come un'istruzione \ti{atomica}, il risultato finale sarà deterministicamente \(a = 2\), e la modellazione con un'unica transizione per processo risulterebbe corretta.

  Se, al contrario, il sistema reale implementa l'istruzione tramite l'esecuzione sequenziale di tre operazioni a livello macchina, la presunta atomicità viene meno:
  \begin{itemize}
    \item \textbf{P1:} \code{load R1, a} \quad \(\rightarrow\) \quad \code{inc R1} \quad \(\rightarrow\) \quad \code{store R1, a}
    \item \textbf{P2:} \code{load R2, a} \quad \(\rightarrow\) \quad \code{inc R2} \quad \(\rightarrow\) \quad \code{store R2, a}
  \end{itemize}
  A causa dell'interleaving in un sistema concorrente, potrebbe verificarsi la seguente esecuzione: P1 carica il valore \(a=0\) in \code{R1}; P2 interrompe P1, carica anch'esso \(a=0\) in \code{R2}, lo incrementa e salva \(a=1\); P1 riprende l'esecuzione, incrementa \code{R1} e salva a sua volta \(a=1\). Il risultato finale è \(1\). Modellare l'incremento come una singola transizione atomica maschererebbe questo comportamento, producendo una verifica formalmente invalida rispetto al sistema reale.
\end{example}

Per illustrare ulteriormente la criticità dell'atomicità, consideriamo una variante del protocollo di mutua esclusione.
In questo protocollo, un processo verifica che la sezione critica sia libera ispezionando il flag dell'avversario e, non appena la trova libera, imposta il proprio flag per indicare la sua presenza, per poi entrare.

La variante del protocollo di mutua esclusione è la seguente:

\begin{figure}{H}
  \centering
  \includegraphics[width=\textwidth]{_images/variant_mutual_exclusion_protocol_code_example}
  \caption{Variante del protocollo di mutua esclusione}\label{fig:variant_mutual_exclusion_protocol_example}
\end{figure}

La correttezza di questo protocollo dipende strettamente da come viene modellata la fase di valutazione e assegnamento.

Possiamo formalizzare il blocco di \code{entry\_protocol} in due modi differenti, ottenendo risultati diametralmente opposti:
\begin{enumerate}
  \item \textbf{Modellazione Atomica (Test-and-Set):}
        Se il sistema dispone di un'istruzione hardware dedicata (come il \keyword{Test-and-Set}), il controllo della guardia e l'assegnamento avvengono in un singolo istante di tempo. Nel Program Graph, ciò si traduce in un'unica transizione dallo stato \(wt_a\) (waiting) allo stato \(cr_a\) (critical), con guardia \(\neg\code{flag}_2\) ed effetto \(\code{flag}_1 := \code{True}\). In questo scenario, il modello è esente da deadlock e garantisce rigorosamente la mutua esclusione.

        \begin{figure}{H}
          \centering
          \includegraphics[width=\textwidth]{_images/automata_with_atomic_transaction_example}
          \caption{Automa per il processo A con transizione atomica}\label{fig:automata_with_atomic_transaction}
        \end{figure}
  \item \textbf{Modellazione Non Atomica:}
        Se il sistema concreto valuta la condizione e aggiorna la variabile in due istruzioni separate, la modellazione deve riflettere questa divisione introducendo uno stato intermedio \(en_a\) (entry). La prima transizione verifica la condizione (\(\neg\code{flag}_2\)) e porta il sistema in \(en_a\); la seconda esegue incondizionatamente l'assegnamento (\(\code{flag}_1 := \code{True}\)) ed entra in \(cr_a\).
        \begin{figure}{H}
          \centering
          \includegraphics[width=\textwidth]{_images/automata_with_non_atomic_transaction_example}
          \caption{Automa per il processo A con transizione non atomica}\label{fig:automata_with_non_atomic_transaction}
        \end{figure}
\end{enumerate}

Nel secondo scenario, la disaggregazione delle azioni genera una vulnerabilità fatale: la \textbf{violazione della mutua esclusione}. È infatti possibile il seguente interleaving (osservabile nel Transition System composto):
\begin{itemize}
  \item Il Processo A esegue il test, trova \(\code{flag}_2 = \code{False}\) e transita nello stato \(en_a\).
  \item Il Processo B viene schedulato \ti{prima} che A imposti il proprio flag. B valuta \(\code{flag}_1 = \code{False}\), transita in \(en_b\).
  \item A questo punto, entrambi i processi si trovano nei rispettivi stati intermedi (\(en_a, en_b\)) ed entreranno immancabilmente in sezione critica (\(cr_a, cr_b\)) l'uno dopo l'altro.
\end{itemize}


\section{Comunicazione e Sincronizzazione}

Finora abbiamo analizzato meccanismi di composizione parallela prevalentemente asincroni, in cui la comunicazione tra processi avviene tramite l'utilizzo di variabili condivise. Tuttavia, in contesti esecutivi in cui non esiste una memoria globale (come nei sistemi distribuiti), è necessario modellare la comunicazione in modo differente, tipicamente tramite lo scambio di messaggi.

\subsection{Comunicazione e Sincronizzazione via Handshake}

Se ci limitiamo al formalismo dei \keyword{Labeled Transition System} (LTS), non abbiamo a disposizione un concetto esplicito di area di memoria o di canale fisico di rete. L'unica entità strutturale a nostra disposizione per far interagire i componenti sono le etichette delle transizioni (le azioni).

L'intuizione alla base è definire un insieme di azioni comuni, condivise dai processi coinvolti. Quando due processi condividono un'azione, affinché uno dei due possa eseguirla, è strettamente necessario che anche l'altro sia nel medesimo istante pronto a eseguire la stessa azione. Questo genera un \ti{parallelismo reale} sulle azioni condivise, mentre sulle azioni puramente locali, o private, continua a vigere il regime di \keyword{interleaving}. Tale modello formale prende il nome di sincronizzazione via \keyword{handshake}.

\begin{definition}{Sincronizzazione via Handshake}
  Siano \(TS_1\) e \(TS_2\) due Transition System, tali che \(TS_i = \langle S_i, A_i, R_i, S_{i0} \rangle\) per \(i \in \{1, 2\}\).
  Definiamo \(H = A_1 \cap A_2\) l'insieme delle \keyword{azioni di sincronizzazione}.

  Il Transition System composto tramite handshake, denotato come \(TS = TS_1 \parallel_H TS_2 = \langle S, A, R, S_0 \rangle\), è così definito:
  \begin{itemize}
    \item \(S = S_1 \times S_2\) \ti{(Stati globali)}
    \item \(A = A_1 \cup A_2\) \ti{(Azioni globali)}
    \item \(S_0 = S_{10} \times S_{20}\) \ti{(Stati iniziali globali)}
  \end{itemize}

  La relazione di transizione \(R\) contiene la transizione \(\langle s_1, s_2 \rangle \xrightarrow{a} \langle s_1', s_2' \rangle\) se e solo se si verifica una delle seguenti due condizioni:
  \begin{enumerate}
    \item \textbf{Esecuzione Sincrona:} \(a \in H\), e per entrambi i processi vale \(s_1 \xrightarrow{a}_1 s_1'\) e \(s_2 \xrightarrow{a}_2 s_2'\). I due processi evolvono simultaneamente eseguendo l'azione condivisa.
    \item \textbf{Interleaving (Azione Locale):} \(a \in A_i \setminus H\), il processo \(i\) evolve localmente \(s_i \xrightarrow{a}_i s_i'\) mentre l'altro processo resta inattivo nel suo stato corrente, ovvero \(s_{3-i}' = s_{3-i}\) (con \(i \in \{1, 2\}\)).
  \end{enumerate}
\end{definition}

\begin{observation}{Modello Ibrido}
  La composizione via handshake rappresenta formalmente un modello ibrido. Costituisce un punto di incontro tra il prodotto puramente sincrono (dove tutto evolve in parallelo) e il prodotto puramente asincrono (dove si applica un interleaving stretto), dove il vincolo sincrono si applica solo sulle azioni di sincronizzazioni \(H\).
\end{observation}

Utilizzando la composizione via handshake, è possibile modellare un protocollo di mutua esclusione che non fa uso di variabili condivise, ma delega il controllo a un processo esterno: un \keyword{Arbitro}.

Consideriamo due processi identici, \(TS_1\) e \(TS_2\), che competono per l'accesso a una risorsa. Ogni processo possiede due stati: \(NC\) (Non Critico) e \(CR\) (Critico). Le transizioni tra questi stati sono etichettate con due azioni: \code{req} (richiesta di accesso) e \code{rel} (rilascio della risorsa).

Se calcolassimo il prodotto puramente asincrono (interleaving puro) tra i due processi (\(TS_1 \parallel TS_2\)), il sistema globale permetterebbe la transizione verso lo stato di errore \(\langle CR_1, CR_2 \rangle\), poiché nessuna entità impedisce ai processi di eseguire l'azione \code{req} in sequenza.

\begin{figure}[H]
  \centering
  \includegraphics[width=\textwidth]{_images/mutual_exclusion_protocol_without_arbiter}
  \caption{Esempio di mutua esclusione senza arbitro}\label{fig:mutual_exclusion_protocol_without_arbiter}
\end{figure}

Per risolvere il problema, introduciamo un terzo automa, l'Arbitro (\(ARB\)), dotato di due stati: \(F\) (Free, risorsa libera) e \(L\) (Locked, risorsa occupata).

Componendo il sistema tramite handshake, imponiamo che le azioni \code{req} e \code{rel} siano \keyword{azioni di sincronizzazione} condivise tra l'interleaving dei due processi e l'arbitro. Formalmente:
\[TS_{globale} = (TS_1 \parallel TS_2) \parallel_{H} ARB \quad \text{con} \quad H = \{\text{req}, \text{rel}\}\]

\begin{figure}[H]
  \centering
  \includegraphics[width=\textwidth]{_images/mutual_exclusion_protocol_with_arbiter}
  \caption{Esempio di mutua esclusione con arbitro}\label{fig:mutual_exclusion_protocol_with_arbiter}
\end{figure}

Quando \(TS_1\) esegue \code{req}, forza l'Arbitro a sincronizzarsi e a spostarsi nello stato \(L\). Finché l'Arbitro si trova in \(L\), non possiede transizioni uscenti etichettate con \code{req}. Di conseguenza, se \(TS_2\) tentasse di eseguire \code{req}, rimarrebbe bloccato, poiché l'azione richiede la partecipazione congiunta dell'Arbitro. L'unica azione accettata dall'Arbitro in \(L\) è \code{rel} (da parte di \(TS_1\)). In questo modo, lo stato invalido \(\langle CR_1, CR_2 \rangle\) diviene irraggiungibile e la mutua esclusione è garantita formalmente.

\begin{note}{}
  L'operatore di composizione parallela \(\parallel_H\) è parametrico rispetto all'insieme di azioni di sincronizzazione \(H\). Questo generalizza i concetti visti in precedenza:
  \begin{itemize}
    \item Se \(H = \emptyset\), l'operatore decade nel \keyword{prodotto interleaving} puro (composizione totalmente asincrona).
    \item Se \(H\) contiene tutte le azioni dell'alfabeto, l'operatore si comporta come un \keyword{prodotto sincrono} puro.
  \end{itemize}
\end{note}

% I'm here
\subsection{Comunicazione via Scambio di Messaggi}

Finora abbiamo applicato la comunicazione via handshake unicamente ai Transition System. Per modellare sistemi software più vicini al codice reale, in cui i processi possiedono variabili ma non condividono uno spazio di memoria (come nei sistemi distribuiti), generalizziamo il concetto ai Program Graph introducendo i \keyword{Channel Systems} (Sistemi a Canali).

Come semplificazione asssumeremo che per questo modello di comunicazione via canali le variabili siano locali e la comunicazione avvenga attraverso \keyword{canali} condivisi.

\begin{definition}{Canali e Capacità}
  Un \keyword{canale} \(c\) è modellato logicamente come una struttura dati di tipo coda FIFO (First-In, First-Out). Ad ogni canale \(c\) sono associati:
  \begin{itemize}
    \item Un dominio \(\text{dom}(c)\), che definisce il tipo di dati dei messaggi trasportabili.
    \item Una capacità \(\text{cap}(c) \ge 0\), che indica il numero massimo di messaggi che il canale può memorizzare simultaneamente.
  \end{itemize}

  Il valore della capacità determina la semantica della comunicazione:
  \begin{itemize}
    \item \textbf{Sincrona:} Se \(\text{cap}(c) = 0\), il canale non possiede memoria. La comunicazione degenera in un handshake stretto: il mittente e il destinatario devono partecipare contemporaneamente alla comunicazione.
    \item \textbf{Asincrona:} Se \(\text{cap}(c) > 0\), il canale agisce da buffer. I processi possono scrivere nel canale in modo asincrono, a patto che questo non sia pieno, e leggere in modo asincrono, a patto che non sia vuoto.
  \end{itemize}
\end{definition}

L'interazione con i canali avviene attraverso primitive specifiche.

\begin{definition}{Azioni di Comunicazione}
  Dato un insieme di variabili \(Var\) e un insieme di canali \(Chan\), l'insieme delle azioni di comunicazione \(Com\) è definito come:
  \[Com = \{c!d, c?x \mid c \in Chan, d \in D, x \in Var \text{ e } \text{dom}(c) \subseteq \text{dom}(x)\}\]
  dove:
  \begin{itemize}
    \item \(c!d\) denota l'operazione di \keyword{scrittura} (send): invio del valore \(d\) sul canale \(c\). L'operazione può essere generalizzata all'invio di un'espressione arbitraria \(\expr \in \Expr\), la quale viene valutata nello stato corrente \(\eta\) prima dell'invio.
    \item \(c?x\) denota l'operazione di \keyword{lettura} (receive): prelievo del primo messaggio disponibile dal canale \(c\) e assegnamento alla variabile locale \(x\).
  \end{itemize}
\end{definition}

\begin{note}{}
  Affinché le operazioni siano ben definite dal punto di vista semantico e dei tipi, deve valere il seguente vincolo sui domini: il dominio del canale deve essere compatibile e incluso nel dominio della variabile di destinazione, ovvero \(\text{dom}(c) \subseteq \text{dom}(x)\), e ovviamente il messaggio inviato deve rispettare il tipo del canale (\(d \in \text{dom}(c)\)).
\end{note}

Definiamo adesso il sistema di canali.

\begin{definition}{Sistema di Canali}
  Un \keyword{Program Graph esteso} su \((Var, Chan)\) è definito come \(PG_i = \langle Loc_i, Act_i, Effect_i, \hookrightarrow_i, Loc_{i0}, g_{i0} \rangle\), dove la relazione di transizione include ora anche le azioni di comunicazione:
  \[\hookrightarrow_i \subseteq Loc_i \times \text{Cond}(Var_i) \times (Act_i \cup Com) \times Loc_i\]

  Un \keyword{Channel System} \(CS\) su \((Var, Chan)\) è la composizione di più Program Graph operanti in parallelo:
  \[CS = PG_1 \parallel PG_2 \parallel \dots \parallel PG_n\]
\end{definition}


Per definire l'LTS risultante, dobbiamo definire:

\begin{itemize}
  \item l'effetto delle azioni nel sistema di comunicazione.
  \item lo spazio degli stati.
\end{itemize}

\begin{definition}{Effetto delle azioni di comunicazione}
  Sia \(c\) un canale. La semantica delle azioni si biforca in base a \(\text{cap}(c)\):

  \textbf{1. Comunicazione Sincrona (\(\text{cap}(c) = 0\)):}
  \begin{itemize}
    \item La transizione \(l_i \xrightarrow{c!d}_i l_i'\) di un processo \(P_i\) è eseguibile \ti{solo se} esiste simultaneamente una transizione \(l_j \xrightarrow{c?x}_j l_j'\) eseguibile da un altro processo \(P_j\).
    \item Le due transizioni avvengono in modo congiunto e l'effetto è l'assegnamento diretto \(x = d\).
  \end{itemize}

  \textbf{2. Comunicazione Asincrona (\(\text{cap}(c) > 0\)):}
  \begin{itemize}
    \item \textbf{Invio (\(c!d\)):} È presente una \keyword{guardia implicita} che richiede che il canale non sia pieno. L'effetto consiste nell'accodare (append) il valore in fondo al buffer FIFO\@:
          \[c = \langle m_1, \dots, m_k \rangle \implies c = \langle m_1, \dots, m_k, d \rangle\]
    \item \textbf{Ricezione (\(c?x\)):} È presente una \keyword{guardia implicita} che richiede che il canale non sia vuoto. L'effetto consiste nel prelevare il messaggio in testa alla coda e assegnarlo alla variabile:
          \[c = \langle m_1, \dots, m_k \rangle \implies c = \langle m_2, \dots, m_k \rangle \text{ e } x = m_1\]
  \end{itemize}
\end{definition}

Per lo spazio degli stati invece è necessario generalizzare la nozione di stato globale introdotta per i semplici Program Graph. In precedenza, uno stato era una coppia formata dalle locazioni e dalla valutazione delle variabili. Ora, poiché i canali possiedono memoria e si comportano come particolari entità di stato condivise, dobbiamo includerne il contenuto.

\begin{definition}{Stato Globale di un Channel System}
  Nel Transition System \(TS(CS)\) risultante dalla composizione, ogni stato è rappresentato da una tupla della forma:
  \[\langle l_1, \dots, l_n, \eta, \chi \rangle\]
  dove:
  \begin{itemize}
    \item \(l_i \in Loc_i\) rappresenta la locazione corrente per ciascun processo \(P_i\).
    \item \(\eta \in \Eval(\Var)\) è la funzione di valutazione delle variabili locali (\(\eta: \Var \to D\)).
    \item \(\chi \in \Eval(Chan)\) è la \keyword{funzione di valutazione dei canali}, che assegna ad ogni canale con \(\text{cap}(c) > 0\) il suo contenuto. Dunque \(\chi\mapsto\text{dom}(c)^*\in \Eval(Chan)\)
  \end{itemize}
\end{definition}

Dunque adesso possiamo formalizzare il transition system per il channel system:

\begin{definition}{Formalizzazione di \(TS(CS)\)}
  Sia \(CS = PG_1 \parallel \dots \parallel PG_n\) un Channel System. Il \keyword{Labeled Transition System} risultante \(TS(CS) = \langle S, Act, \to, S_0, AP, L \rangle\) è formalmente definito come segue:

  \begin{itemize}
    \item \textbf{Stati (\(S\)):} L'insieme degli stati è \(S = (Loc_1 \times \cdots \times Loc_n) \times \Eval(\Var) \times \Eval(Chan)\).
    \item \textbf{Azioni (\(Act\)):} L'insieme delle azioni globali è \(Act = \biguplus_{1 \le i \le n} Act_i \uplus \{\tau\}\), dove \(\tau\) è un simbolo di \keyword{azione interna} che viene usata per modellare la sincronizzazione.
    \item \textbf{Stati Iniziali (\(S_0\)):} Uno stato \(\langle l_1, \dots, l_n, \eta, \chi \rangle\) appartiene a \(S_0\) se:
          \begin{itemize}
            \item \(l_i \in Loc_{i0}\) per ogni \(1 \le i \le n\).
            \item \(\eta \models \bigwedge_{i=1}^n g_{i0}\) (la valutazione iniziale \(\eta\) soddisfa la congiunzione delle guardie iniziali di tutti i processi).
            \item \(\chi = \chi_0\), ovvero \(\chi(c) = \epsilon\) per ogni canale \(c \in Chan\) (i canali sono inizialmente vuoti).
          \end{itemize}
    \item \textbf{Proposizioni Atomiche e Labelling (\(AP, L\)):} L'insieme \(AP\) e la funzione di etichettatura \(L\) rimangono invariati rispetto ai Program Graph standard, includendo locazioni e condizioni sulle variabili.
    \item \textbf{Relazione di transizione (\(\to\)):} definita come segue.
  \end{itemize}
\end{definition}

La relazione di transizione globale \(\to\) descrive come si aggiorna lo stato globale a seguito dell'esecuzione di azioni private o di interazione con i canali. Possiamo suddividere la sua formalizzazione in tre casi distinti.

\begin{itemize}
  \item \textbf{Azioni Locali:} Se il processo \(P_i\) esegue un'azione locale standard \(a \in Act_i\), il suo stato avanza e si applicano gli effetti sulle variabili, senza toccare i canali.
        \[ \text{Se } l_i \xrightarrow{g:a}_i l_i' \text{ e } \eta \models g \]
        \[ \implies \langle l_1, \dots, l_i, \dots, l_n, \eta, \chi \rangle \xrightarrow{a} \langle l_1, \dots, l_i', \dots, l_n, \eta', \chi \rangle \]
        dove \(\eta' = \text{Effect}(a, \eta)\).
  \item \textbf{Comunicazione Asincrona:} Le azioni di I/O asincrone aggiornano la valutazione dei canali \(\chi\) (e possibilmente \(\eta\)) ed etichettano la transizione globale con \(\tau\).
        \begin{itemize}
          \item \textbf{Azione di Scrittura (Send):} Richiede canale non pieno. L'espressione valutata o il dato \(d\) viene accodato.
                \[ \text{Se } l_i \xrightarrow{g:c!d}_i l_i', \quad \eta \models g, \quad \text{len}(\chi(c)) = k < \text{cap}(c), \quad \chi(c) = \langle d_1, \dots, d_k \rangle \]
                \[ \implies \langle l_1, \dots, l_i, \dots, l_n, \eta, \chi \rangle \xrightarrow{\tau} \langle l_1, \dots, l_i', \dots, l_n, \eta, \chi' \rangle \]
                dove \(\chi' = \chi[c := \langle d_1, \dots, d_k, d \rangle]\). La sintassi \(\chi[c := \sigma]\) indica una funzione identica a \(\chi\) ovunque, tranne che in \(c\), cui associa la sequenza \(\sigma\). La funzione di valutazione delle variabili \(\eta\) rimane invariata.

          \item \textbf{Azione di Lettura (Receive):} Richiede canale non vuoto. Si preleva la testa della coda e si aggiorna una variabile locale.
                \[ \text{Se } l_i \xrightarrow{g:c?x}_i l_i', \quad \eta \models g, \quad \text{len}(\chi(c)) = k > 0, \quad \chi(c) = \langle d_1, \dots, d_k \rangle \]
                \[ \implies \langle l_1, \dots, l_i, \dots, l_n, \eta, \chi \rangle \xrightarrow{\tau} \langle l_1, \dots, l_i', \dots, l_n, \eta', \chi' \rangle \]
                dove \(\chi' = \chi[c := \langle d_2, \dots, d_k \rangle]\) (avanzamento della coda) e \(\eta' = \eta[x := d_1]\) (assegnamento del valore estratto alla variabile \(x\)).
        \end{itemize}
  \item \textbf{Comunicazione Sincrona:} L'handshake richiede la presenza simultanea di una scrittura e di una lettura da parte di due processi distinti \(P_i\) e \(P_j\) sullo stesso canale \(c\). Il dato transita in modo sincrono.

        \[ \text{Se } l_i \xrightarrow{g_1:c!d}_i l_i' \text{ e } l_j \xrightarrow{g_2:c?x}_j l_j', \quad \eta \models g_1 \land g_2, \quad \text{cap}(c) = 0, \quad i \neq j \]
        \[ \implies \langle l_1, \dots, l_i, \dots, l_j, \dots, l_n, \eta, \chi \rangle \xrightarrow{\tau} \langle l_1, \dots, l_i', \dots, l_j', \dots, l_n, \eta', \chi \rangle \]
        dove \(\eta' = \eta[x := d]\). L'azione globale è ancora \(\tau\). La valutazione dei canali \(\chi\) rimane intoccata, poiché il canale non conserva mai informazione.
\end{itemize}

\begin{observation}{Approcci alla Composizione}
  Esiste una differenza strutturale nel modo in cui abbiamo derivato il Transition System composto per i sistemi con memoria condivisa rispetto ai sistemi con comunicazione a canali:
  \begin{itemize}
    \item \textbf{Modello a Memoria Condivisa:} L'approccio classico procede per gradi. Prima si esegue la composizione parallela sintattica tra i Program Graph (\(PG_{comp} = PG_1 \parallel \dots \parallel PG_n\)), sfruttando l'interleaving e la gestione condivisa delle variabili. Successivamente, l'unico Program Graph risultante viene \q{srotolato} (unfolded) nel suo LTS corrispondente \(TS(PG_{comp})\).
    \item \textbf{Modello a Canali (Channel System):} A causa della natura dei canali (che agiscono come code FIFO indipendenti) e delle primitive di I/O asincrone/sincrone, non definiamo un operatore di composizione nativo a livello di Program Graph. L'approccio adottato consiste nel mappare \ti{direttamente} l'insieme dei Program Graph separati, unito all'insieme dei canali, in un unico e grande Transition System globale \(TS(CS)\), applicando una funzione di traduzione diretta che incorpora la semantica dei canali nella definizione della relazione di transizione.
  \end{itemize}
  Se il sistema da modellare è \keyword{ibrido} (presenza simultanea di variabili condivise e canali di comunicazione), non si può procedere in due fasi separate. La traduzione in LTS deve necessariamente fondere i due approcci in un'unica definizione monolitica della relazione di transizione globale.
\end{observation}



\section{Verso i Modelli a Tempo Continuo}

La modellazione vista finora si è concentrata su sistemi discreti, in cui il tempo avanza concettualmente a \q{scatti} in corrispondenza dell'esecuzione delle transizioni atomiche. Esiste tuttavia un'ampia classe di sistemi la cui correttezza dipende crucialmente dallo scorrere del \keyword{tempo reale} e dall'evoluzione di grandezze fisiche continue.

L'idea fondamentale per modellare tali sistemi è iniettare all'interno del formalismo discreto delle variabili che non evolvono in base ad assegnamenti discreti operati dal software, ma in base a leggi matematiche che dipendono dal tempo (es.\ temperatura, velocità, pressione).

\begin{definition}{Variabile Tempo Reale}
  Introduciamo una singola variabile continua per rappresentare il tempo assoluto, che denoteremo tipicamente con \(t\). A differenza delle variabili di programma classiche, il cui dominio è spesso finito o numerabile, la variabile tempo varia nel dominio dei numeri reali non negativi:
  \[ \text{dom}(t) = \R_{\ge 0} \]
  Questa variabile evolve autonomamente con derivata unitaria (\(\dot{t} = 1\)), misurando il tempo reale trascorso.
\end{definition}

\begin{observation}{Esplosione dello Spazio degli Stati e Decidibilità}
  L'introduzione di una variabile a dominio continuo altera drasticamente la trattabilità matematica e computazionale del modello.

  Ricordiamo che in un Transition System generato da variabili, lo spazio degli stati \(S\) è dato (o contiene) l'insieme di tutte le possibili valutazioni delle variabili. Una funzione di valutazione \(\eta\) mappa l'insieme delle variabili \(\Var\) nel loro dominio \(D\):
  \[ \eta: \Var \to D \]
  Se l'unica variabile di interesse è il tempo (\(\Var = \{t\}\)) e il dominio è \(D = \R_{\ge 0}\), il numero totale di valutazioni distinte corrisponde al numero di funzioni possibili:
  \[ \abs{D}^{\abs{\Var}} = \abs{\R_{\ge 0}}^1 = \abs{\R_{\ge 0}} \]
  Questo significa che lo spazio degli stati del sistema non è più finito, né numerabilmente infinito (come \(\N\)), ma possiede la cardinalità del continuo (non numerabile).

  Dal punto di vista pratico dell'Automated Software Verification, questo genera una barriera computazionale enorme. Gli algoritmi classici di verifica (Model Checking), che si basano sull'esplorazione esaustiva degli stati raggiungibili, falliscono in presenza di uno spazio non numerabile. Il problema della verifica per modelli con variabili continue diviene, in linea di principio, \keyword{indecidibile}. Sarà pertanto necessario introdurre tecniche formali di \keyword{astrazione} avanzate per ricondurre questo spazio continuo a un sistema finito equivalente ai fini della verifica, tema che verrà esplorato nello studio dei \keyword{Timed Automata}.
\end{observation}